A key roadblock to applying supervised learning to real applications is the need to assemble a large-enough labeled dataset for the intended task.
Modern deep learning pipelines are especially data-hungry.
In many cases, the acquisition of a large dataset of \emph{unlabeled} features is rather affordable.
However, providing reliable labels for each example is cost-prohibitive, often requiring expensive, time-consuming work from human experts.
This tradeoff is especially apt in our motivating application: classifying medical images where
images are collected in the course of routine care and easily available by querying a hospital's electronic records.
However, labeling images often requires clinical staff with years of training to spend minutes per image.

If only a tiny labeled set is available but we can access a big \emph{unlabeled} set of images, one promising approach is \emph{semi-supervised learning} (SSL)~\citep{zhuSemiSupervisedLearningLiterature2005,vanengelenSurveySemisupervisedLearning2020}. 
Recent years have seen amazing progress on standard benchmarks such as recognizing address digits from photos of houses (SVHN;~\citet{netzerReadingDigitsNatural2011}).
With only 100 labeled examples per digit class, a supervised neural net's error rate is roughly 12\%. 
Using a large unlabeled set, the FixMatch SSL method~\citep{sohnFixMatchSimplifyingSemisupervised2020} delivers error below 2.5\%, while even more recent work has pushed below 2\%~\citep{xuDPSSLRobustSemisupervised2021, hanUnsupervisedSemanticAggregation2020}.

Unfortunately, common benchmarks like SVHN may be too optimistic. 
In real tasks, unlabeled sets will be collected automatically at scale for convenience, and thus \emph{uncurated}: they may differ from the labeled set in terms of represented classes, class frequencies, or even features.
% To maintain the efficiency that motivates SSL, the effort required to curate the unlabeled set -- that is, apply labels and assess differences -- remains prohibitive.
Effective SSL must improve accuracy despite such uncurated data.

Off-the-shelf SSL using mismatched unlabeled sets often predicts \emph{worse} than just ignoring unlabeled data~\citep{oliverRealisticEvaluationDeep2018,calderon-ramirezDealingScarceLabelled2021}.
Recent methods try to be robust to unlabeled sets that differ from the labeled set (see Tab.~\ref{tab:related_work_comparison}).
% line of work called ``safe'' or ``open-set'' SSL.
%deliver accuracy that is at least as good as using the labeled set alone but better when possible.
The dominant paradigm is intuitive: 
identify examples in the unlabeled set that are \emph{out-of-distribution} (OOD), then remove or downweight them \citep{calderon2022semi,chen2022semi, he2022not}.
We find this line of work delivers insufficient gains in accuracy, while adding complexity due to OOD detection and discarding a substantial amount of unlabeled data.
%The resulting algorithms thus depend heavily on the OOD detector's performance.
% which adds additional complexity \todo{while delivering only modest gains.} 

% Larger, more reliable gains are needed to enable SSL in real uncurated applications.
%\todo{Further, viewing OOD samples as completely useless limited the possible performance of SSL algorithms in real applications, since a substantial portion of the unlabeled data could be discarded}. 

This study makes 3 contributions toward robust SSL backed by reproducible experiments\footnote{Code and Heart2Heart data: \href{https://github.com/tufts-ml/fix-a-step}{github.com/tufts-ml/fix-a-step}}.
First, we challenge the dominant paradigm of filtering out or downweighting OOD examples in the unlabeled set.
Our experiments suggest that even perfect OOD filtering, which is unrealistic in practice, does not perform well (see Fig.~\ref{fig:results-cifar10}).
Instead of viewing OOD images as probably harmful, we argue for a \textbf{new paradigm: OOD images from uncurated unlabeled sets are possibly helpful}.
%this sentence need to break down?

Second, following this paradigm we introduce \textbf{a new training procedure called \emph{Fix-A-Step} that delivers accuracy gains from uncurated unlabeled sets}.
When applied to repair several deep SSL methods across a range of labeled-unlabeled class mismatch levels, our Fix-A-Step improves predictions better than alternative methods while being substantially simpler and faster too.

Finally, we offer a \textbf{new SSL benchmark called Heart2Heart that uses truly uncurated unlabeled medical images and assesses cross-hospital generalization}.
%For too long, SSL methods advances have been evaluated only on the same repurposed generic object recognition datasets. 
Using three inter-operable open-access datasets -- TMED~\citep{huangNewSemisupervisedLearning2021}, CAMUS~\citep{leclercDeepLearningSegmentation2019}, and Unity~\citep{howardAutomatedLeftVentricular2021} --
we pursue a clinically-relevant problem: recognizing the view type of an echocardiogram image of the heart.
Future methods that learn from limited data can follow our reproducible protocol. 
We hope this new Heart2Heart benchmark enables authentic SSL applications in medicine and ultimately improves care for patients with heart disease.
